\section{Introduction}
\label{sec:intro}

Identifying which function is responsible for a reported bug---fault
localisation---is one of the most time-consuming steps in software
maintenance~\cite{wong16survey}.
Traditional automated approaches such as spectrum-based FL (SBFL) rank
functions by test-coverage patterns~\cite{jones05tarantula}, requiring both a
test suite and the ability to execute the program.
Large language models (LLMs) offer a complementary capability: they can reason
over natural-language bug descriptions and source code simultaneously, and can
be equipped with \emph{tools} that retrieve task-relevant information
on-the-fly~\cite{yao23react}.
However, LLMs cannot read an entire production codebase within their context
window, so retrieval quality is a key bottleneck.

A code call-graph provides a structured, navigable index of which functions
call which others.
Given a call-graph, a tool-using agent can iteratively narrow the search---
following callers upward, callees downward, and taint paths through the
data-flow---without loading the full source.
Alternatively, a breadth-first slice of the same graph can be serialised as
plain text and injected into the prompt, trading repeated tool-call overhead
for a single, larger context window.
Both strategies have been proposed in recent work, but a direct controlled
comparison on a real, well-known codebase under identical conditions is lacking.

We build two call-graphs for the \sqlite~3.30.1 source: a compile-time graph
extracted with \clang~18.1.3 (417 functions, 2{,}438 call edges, 84 taint
edges) and a syntactic graph extracted with tree-sitter/graphify.
We then run a controlled experiment across five conditions (no-context,
tool-agent on each graph, text-slice on each graph) and four tool-ablation
variants, using two instruction-tuned LLMs (\qwen and \nemotron) on six
fault-localisation benchmark questions.

\paragraph{Research questions.}
\begin{description}
  \item[\textbf{RQ1}] Does call-graph context (tool or text) improve
    fault-localisation accuracy over a no-context baseline?
  \item[\textbf{RQ2}] Does tool-based graph navigation yield higher accuracy
    than delivering the same graph as a text slice?
  \item[\textbf{RQ3}] Which graph-navigation tools are necessary for accuracy?
  \item[\textbf{RQ4}] How often do graph-navigating LLMs cite call edges that
    the graph does not contain?
\end{description}

\paragraph{Main findings.}
\begin{enumerate}
  \item Graph context provides only marginal accuracy gains: neither model
    exceeds one correct answer out of five (Qwen, excluding Q2) or six across any single condition.
    Furthermore, 15 of 48 answer slots (31\,\%) arise from function-pointer gaps that make
    the target structurally unreachable in both graphs.
  \item Text slices match or exceed the tool condition in exact accuracy at
    about ten to twelve times lower token cost for \qwen and sixteen to twenty-three
    times lower for \nemotron, though the exact pattern is model-dependent:
    \nemotron solved Q5 in both tool conditions while \qwen succeeded only in the text conditions.
  \item Source-only retrieval (no edge-traversal tools) consumes sixty times more tokens
    than a text slice with no accuracy gain; caller/callee tools were associated with
    a 2.4-fold lower rate of edges not supported by the graph (12\,\% vs.\ 28\,\%).
\end{enumerate}

\paragraph{Contributions.}
\begin{enumerate}
  \item Two open call-graph databases for \sqlite~3.30.1 (compile-time and
    syntactic), together with graph query tool suites (7 tools for CT, 5 for SYN),
    a text-slice generator, and a resumable multi-condition evaluation harness.
  \item An empirical study comparing tool-based and text-slice context delivery
    across two LLMs and nine experimental conditions, with a detailed error
    classification and path-correctness analysis.
  \item An annotated dataset of 48 answer slots (including 43 failure slots) and 133
    edge-verification records identifying when graph-conditioned LLMs cite edges
    absent from the graph.
\end{enumerate}
